%!TEX root = ../template.tex
%%%%%%%%%%%%%%%%%%%%%%%%%%%%%%%%%%%%%%%%%%%%%%%%%%%%%%%%%%%%%%%%%%%%
%% abstract-pt.tex
%% NOVA thesis document file
%%
%% Abstract in Portuguese
%%%%%%%%%%%%%%%%%%%%%%%%%%%%%%%%%%%%%%%%%%%%%%%%%%%%%%%%%%%%%%%%%%%%

\typeout{NT FILE abstract-pt.tex}%
Os drones tornaram-se cada vez mais populares nos últimos anos devido à sua versatilidade em múltiplas ocasiões, seja no contexto militar, civil, ou para transporte, mapeamento ou vigilância. Com este aumento, não só têm surgido soluções novas e inovadoras para diversos problemas, como também atores mal-intencionados têm tirado proveito destas mesma tecnologias.
Esta tese aborda este problema através do estudo de jogos de perseguição-evasão, com foco no planeamento de trajetórias para garantir um percurso livre de colisões numa perseguição envolvendo múltiplos drones a um infrator. O rastreamento de referência foi realizado através de um controlador Millinger já validado, sendo a parte principal da tese baseada em funções de Lyapunov de controlo numa dada manifold.

\keywords{
  \gls{CLF} \and
  UAV \and
  Path planning
}